\documentclass[11pt,a4paper]{article}
\usepackage{amsmath,amssymb,amsthm}
\newtheorem{theorem}{Theorem}
\begin{document}
\section{Euler's formula}
For a connected planar graph with $V$ vertices, $E$ edges and $F$ faces,
\begin{equation}\label{eq:euler}
V-E+F=2.
\end{equation}

\begin{theorem}\label{thm:k5}
The complete graph $K_5$ is not planar.
\end{theorem}
\begin{proof}
$K_5$ has $V=5$ and $E=10$. If it were planar, \eqref{eq:euler} would give $F=7$. Each face is bounded by at least three edges, and each edge borders at most two faces, so $3F\le 2E$, that is $21\le 20$, a contradiction.
\end{proof}

\section{Degrees}
The handshake lemma states
\[
\sum_{v}\deg(v)=2E.
\]
Consequently the number of vertices of odd degree is even.

\subsection{Trees}
A tree on $n$ vertices has $n-1$ edges. Conversely, a connected graph with $n$ vertices and $n-1$ edges is a tree.

Theorem~\ref{thm:k5} and the handshake lemma are both consequences of counting; see equation~\eqref{eq:euler} on page~\pageref{eq:euler}.
\end{document}
